\section{A restricted first stage and an unrestricted upper boundary}\label{sec:ramification}
We now state the inherited arithmetic contract used by the final affine reduction. These are actual geometry inputs, not consequences of a finite matrix scan.
\begin{description}
\item[H1.] The actual primitive Frey family, its specified twist and integral $\Oloc$ action, determinant $\chi$, and natural depth one endpoint conventions are fixed.
\item[H2.] At every variable bad place the actual inertia has $I+N$ with $N^2=0$ and $N$ in $7M_2(\Oloc)$. Its primitive line factorization and same frame Frobenius obey the toric equations. Good inertia is trivial.
\item[H3.] Outside the four fixed rational primes $2,3,5,7$, every possible bad place divides $C$. The reference is unramified there.
\item[H4.] At each fixed place over $2,3,5,7$, the actual rank two trace is $\chi+1$, and a local element has $\chi-1$ a unit modulo seven.
\item[H5.] The four frozen beta supports, its stated all inert descent, and local natural class zero at every beta support place are retained.
\item[H6.] The complete level seven period criterion and natural node thickness are used for the local Frey positive control only.
\end{description}
The exact original contracts remain in the source archive. Preservation of their bytes is not independent mathematical acceptance of the geometric chain.

\begin{proposition}[No auxiliary diagonal inertia]\label{prop:pinertia}
Under H1--H3, the first cochain $p$ is unramified outside the fixed primes and the complete beta ramification support. It factors as a cochain through the corresponding restricted Galois group.
\end{proposition}
\begin{proof}
At a variable bad place, the reference inertia is identity. Reduce $N/7$ modulo seven. Its traceless matrix is
\[
\begin{pmatrix}p(T)&q(T)\\\beta(T)&-p(T)\end{pmatrix},
\]
and nilpotence gives $p(T)^2+q(T)\beta(T)=0$. If beta is unramified, its value is zero on every inertia element. Thus $p(T)^2=0$ in the field $k$, implying $p(T)=0$. This square zero inference is in a field, not in the ring $\Oloc/49$, which has nilpotents. At good places both representations are unramified, giving the same conclusion.

Let $I$ be an omitted inertia subgroup. The residual classes and $p$ all vanish on $I$. The multiplication law \eqref{eq:p-cocycle} yields $p(gi)=p(g)=p(ig)$ for $i\in I$. It also gives vanishing on conjugates and products of omitted inertia elements. Continuity extends this to their closed normal closure. The cochain descends to the restricted quotient.
\end{proof}
This result permits a finite Stage P computation even if $C$ has many extra divisors unramified in beta. The set of all those divisors need not be bounded. The theorem controls only the diagonal first cochain.

If beta is unramified and $q(T)$ is nonzero, the first digit inertia is purely upper. It can survive in the actual toric local category. The archive retains an exact local Frey positive control at $5167$, with $A=3238$, $C=5167$ and a local fifth root $B$ satisfying the equation. Its leading residue is $3658$, and simple root lifting supplies the root to every precision. This $B$ is a local number, not an ordinary integer solution.

\subsection{Upper values at beta support are fixed after \texorpdfstring{$p$}{p}}
At beta support, the natural class is identically zero locally, and the coefficient action is trivial modulo seven. Relative to the same reference, put $B_0(F)=b_0(F)/7$ modulo seven. The local packet yields
\begin{align}
cp(F)-bp(T)&=-hc,\label{eq:p-support}\\
cq(T)&=-p(T)^2,\label{eq:q-inertia}\\
c(q(F)+B_0(F))&=-p(T)p(F).\label{eq:q-frob}
\end{align}
The reference term in the last equation is retained. It is not reset to zero by an independent local basis choice. Once $p$ and beta are fixed, all upper values on the full beta supporting local group are fixed at this precision. Two upper corrections with the same data therefore have zero local difference there.

At beta unramified bad places, upper inertia may still occur when the local beta class splits. The exact matrices $N=7uE_{12}$ and $F=\left(\begin{smallmatrix}Q&R\\0&1\end{smallmatrix}\right)$ provide all precision toric controls. Hence the attempted theorem confining \emph{every} adjoint ramification to beta support is false. The finite restricted set is justified for Stage P. Stage Q retains its own auxiliary arithmetic problem.
